\begin{abstract}
From Apple Music to Spotify, song suggestion algorithms have become a hallmark of music streaming platforms. Many of these platforms base music recommendations on surface-level factors such as genre or artist, but taking a deeper look at audio data suggests that there may be a way to provide users with a more personalized experience. This project aims to take a closer look at what makes audio files unique, utilizing visual information via spectrograms and sentiment encoding via lyrics. By using a multimodal system and training our model with these factors in addition to song metadata, we aimed to create a more accurate and refined song recommendation system based on existing user preferences.

A link to the code can be found hosted on GitHub at \href{https://github.com/arenvista/AudioRecomendation}{https://github.com/arenvista/AudioRecomendation}.
\end{abstract}